% Windows in time (Section 3.3): where communication happens. One gate axis for three bands.
%   top     cuPauliProp's distributed mode: a global merge after every gate (22 collective calls per gate on
%           every rank, logs/pp_prof4n-59470888.out)
%   middle  SectorQ on four nodes over three windows: inside a window every node works on its own sectors
%           (a sector has the color of its node); at the end of a window the output is already divided into
%           pieces of all next sectors; the exchange at the boundary puts every next sector on its node
%   bottom  measured number of terms over the same three windows (kicked Ising, cutoff 1e-5, windows 8-10,
%           results/profile/ki_lab_1e5.csv, job 59488608): 1.36 billion at the peak, 159 million at the
%           boundary, because a window ends right after a truncation
% generated from results/profile/ki_lab_1e5.csv (kicked Ising, cutoff 1e-5, job 59488608): terms before every
% operation of windows 8-10 (operations 1016-1397), one window per 3.5 cm starting at x = 2.3, 6.2 and 10.1
\def\curvepts{(2.300,-6.096) (2.328,-6.096) (2.355,-6.096) (2.383,-6.096) (2.410,-6.096) (2.438,-6.096) (2.465,-6.096) (2.493,-6.096) (2.520,-6.096) (2.548,-6.096) (2.576,-6.096) (2.603,-6.096) (2.631,-6.096) (2.658,-6.096) (2.686,-6.096) (2.713,-6.096) (2.741,-6.096) (2.769,-6.096) (2.796,-6.096) (2.824,-6.096) (2.851,-6.096) (2.879,-6.096) (2.906,-6.096) (2.934,-6.095) (2.961,-6.090) (2.989,-6.073) (3.017,-6.058) (3.044,-6.051) (3.072,-6.094) (3.099,-6.081) (3.127,-6.073) (3.154,-6.071) (3.182,-6.071) (3.209,-6.071) (3.237,-6.071) (3.265,-6.048) (3.292,-6.018) (3.320,-6.018) (3.347,-6.094) (3.375,-6.094) (3.402,-6.092) (3.430,-6.083) (3.457,-6.051) (3.485,-6.010) (3.513,-5.942) (3.540,-5.854) (3.568,-5.725) (3.595,-5.550) (3.623,-6.093) (3.650,-6.070) (3.678,-6.040) (3.706,-6.019) (3.733,-6.012) (3.761,-6.012) (3.788,-5.983) (3.816,-5.910) (3.843,-5.840) (3.871,-5.838) (3.898,-6.095) (3.926,-6.083) (3.954,-6.062) (3.981,-6.013) (4.009,-5.947) (4.036,-5.850) (4.064,-5.732) (4.091,-5.562) (4.119,-5.366) (4.146,-5.100) (4.174,-6.096) (4.202,-6.063) (4.229,-6.040) (4.257,-6.023) (4.284,-6.019) (4.312,-5.965) (4.339,-5.890) (4.367,-5.879) (4.394,-5.879) (4.422,-5.879) (4.450,-6.098) (4.477,-6.095) (4.505,-6.083) (4.532,-6.062) (4.560,-6.013) (4.587,-5.956) (4.615,-5.895) (4.643,-5.834) (4.670,-5.735) (4.698,-5.640) (4.725,-6.090) (4.753,-6.087) (4.780,-6.085) (4.808,-6.085) (4.835,-6.082) (4.863,-6.074) (4.891,-6.065) (4.918,-6.065) (4.946,-6.065) (4.973,-6.065) (5.001,-6.089) (5.028,-6.089) (5.056,-6.089) (5.083,-6.089) (5.111,-6.089) (5.139,-6.087) (5.166,-6.087) (5.194,-6.087) (5.221,-6.087) (5.249,-6.086) (5.276,-6.090) (5.304,-6.090) (5.331,-6.090) (5.359,-6.090) (5.387,-6.090) (5.414,-6.090) (5.442,-6.090) (5.469,-6.090) (5.497,-6.090) (5.524,-6.090) (5.552,-6.090) (5.580,-6.090) (5.607,-6.090) (5.635,-6.090) (5.662,-6.090) (5.690,-6.090) (5.717,-6.090) (5.745,-6.090) (5.772,-6.090) (5.800,-6.090) (6.200,-6.090) (6.228,-6.090) (6.255,-6.090) (6.283,-6.090) (6.310,-6.090) (6.338,-6.090) (6.365,-6.090) (6.393,-6.090) (6.420,-6.090) (6.448,-6.090) (6.476,-6.090) (6.503,-6.090) (6.531,-6.090) (6.558,-6.090) (6.586,-6.090) (6.613,-6.090) (6.641,-6.089) (6.669,-6.090) (6.696,-6.089) (6.724,-6.089) (6.751,-6.089) (6.779,-6.089) (6.806,-6.089) (6.834,-6.087) (6.861,-6.082) (6.889,-6.065) (6.917,-6.048) (6.944,-6.088) (6.972,-6.084) (6.999,-6.068) (7.027,-6.053) (7.054,-6.049) (7.082,-6.047) (7.109,-6.047) (7.137,-6.047) (7.165,-6.019) (7.192,-5.981) (7.220,-6.089) (7.247,-6.089) (7.275,-6.089) (7.302,-6.086) (7.330,-6.076) (7.357,-6.043) (7.385,-6.001) (7.413,-5.936) (7.440,-5.849) (7.468,-5.726) (7.495,-6.093) (7.523,-6.061) (7.550,-6.033) (7.578,-5.994) (7.606,-5.965) (7.633,-5.955) (7.661,-5.955) (7.688,-5.920) (7.716,-5.826) (7.743,-5.733) (7.771,-6.098) (7.798,-6.094) (7.826,-6.080) (7.854,-6.058) (7.881,-6.009) (7.909,-5.942) (7.936,-5.846) (7.964,-5.734) (7.991,-5.577) (8.019,-5.387) (8.046,-6.103) (8.074,-6.071) (8.102,-6.032) (8.129,-6.004) (8.157,-5.983) (8.184,-5.978) (8.212,-5.914) (8.239,-5.817) (8.267,-5.799) (8.294,-5.799) (8.322,-6.106) (8.350,-6.104) (8.377,-6.100) (8.405,-6.088) (8.432,-6.067) (8.460,-6.022) (8.487,-5.969) (8.515,-5.907) (8.543,-5.849) (8.570,-5.759) (8.598,-6.100) (8.625,-6.091) (8.653,-6.085) (8.680,-6.083) (8.708,-6.083) (8.735,-6.079) (8.763,-6.069) (8.791,-6.055) (8.818,-6.055) (8.846,-6.055) (8.873,-6.099) (8.901,-6.099) (8.928,-6.098) (8.956,-6.098) (8.983,-6.098) (9.011,-6.098) (9.039,-6.096) (9.066,-6.095) (9.094,-6.095) (9.121,-6.094) (9.149,-6.099) (9.176,-6.099) (9.204,-6.099) (9.231,-6.099) (9.259,-6.099) (9.287,-6.099) (9.314,-6.099) (9.342,-6.099) (9.369,-6.099) (9.397,-6.099) (9.424,-6.099) (9.452,-6.099) (9.480,-6.099) (9.507,-6.099) (9.535,-6.099) (9.562,-6.099) (9.590,-6.099) (9.617,-6.099) (9.645,-6.099) (9.672,-6.099) (9.700,-6.099) (10.100,-6.099) (10.128,-6.099) (10.155,-6.099) (10.183,-6.099) (10.210,-6.099) (10.238,-6.099) (10.265,-6.099) (10.293,-6.099) (10.320,-6.099) (10.348,-6.099) (10.376,-6.099) (10.403,-6.099) (10.431,-6.099) (10.458,-6.099) (10.486,-6.099) (10.513,-6.099) (10.541,-6.099) (10.569,-6.099) (10.596,-6.099) (10.624,-6.099) (10.651,-6.099) (10.679,-6.099) (10.706,-6.098) (10.734,-6.097) (10.761,-6.092) (10.789,-6.077) (10.817,-6.098) (10.844,-6.096) (10.872,-6.091) (10.899,-6.076) (10.927,-6.061) (10.954,-6.055) (10.982,-6.053) (11.009,-6.053) (11.037,-6.052) (11.065,-6.026) (11.092,-6.099) (11.120,-6.098) (11.147,-6.098) (11.175,-6.098) (11.202,-6.094) (11.230,-6.085) (11.257,-6.055) (11.285,-6.017) (11.313,-5.959) (11.340,-5.882) (11.368,-6.102) (11.395,-6.071) (11.423,-6.033) (11.450,-6.000) (11.478,-5.955) (11.506,-5.919) (11.533,-5.905) (11.561,-5.905) (11.588,-5.865) (11.616,-5.755) (11.643,-6.107) (11.671,-6.106) (11.698,-6.102) (11.726,-6.089) (11.754,-6.068) (11.781,-6.023) (11.809,-5.963) (11.836,-5.874) (11.864,-5.768) (11.891,-5.619) (11.919,-6.113) (11.946,-6.079) (11.974,-6.040) (12.002,-5.995) (12.029,-5.962) (12.057,-5.938) (12.084,-5.931) (12.112,-5.859) (12.139,-5.745) (12.167,-5.723) (12.194,-6.115) (12.222,-6.115) (12.250,-6.112) (12.277,-6.108) (12.305,-6.096) (12.332,-6.077) (12.360,-6.035) (12.387,-5.987) (12.415,-5.931) (12.443,-5.875) (12.470,-6.111) (12.498,-6.098) (12.525,-6.086) (12.553,-6.079) (12.580,-6.074) (12.608,-6.074) (12.635,-6.070) (12.663,-6.056) (12.691,-6.037) (12.718,-6.037) (12.746,-6.111) (12.773,-6.111) (12.801,-6.111) (12.828,-6.109) (12.856,-6.109) (12.883,-6.109) (12.911,-6.108) (12.939,-6.106) (12.966,-6.104) (12.994,-6.104) (13.021,-6.111) (13.049,-6.110) (13.076,-6.110) (13.104,-6.109) (13.131,-6.109) (13.159,-6.109) (13.187,-6.109) (13.214,-6.109) (13.242,-6.109) (13.269,-6.109) (13.297,-6.111) (13.324,-6.111) (13.352,-6.111) (13.380,-6.111) (13.407,-6.111) (13.435,-6.111) (13.462,-6.111) (13.490,-6.111) (13.517,-6.111) (13.545,-6.111) (13.572,-6.111) (13.600,-6.111)}
\def\truncticks{2.520,2.796,3.072,3.347,3.623,3.898,4.174,4.450,4.725,5.001,5.276,5.552,5.800,6.393,6.669,6.944,7.220,7.495,7.771,8.046,8.322,8.598,8.873,9.149,9.424,9.700,10.265,10.541,10.817,11.092,11.368,11.643,11.919,12.194,12.470,12.746,13.021,13.297,13.572,13.600}

\begin{tikzpicture}[x=1cm, y=1cm, font=\rmfamily\footnotesize]
  \tikzset{sec/.style={draw=black!70, line width=0.3pt, anchor=west, inner sep=0pt, minimum width=0.46cm, minimum height=0.32cm,
                       font=\rmfamily\scriptsize},
           pc/.style={draw=black!70, line width=0.25pt, anchor=west, inner sep=0pt, minimum width=0.11cm, minimum height=0.32cm}}
  % ================= cuPauliProp: a global step after every gate =================
  \node[layer, anchor=north west, minimum width=13.9cm, minimum height=1.12cm] at (0,0) {};
  \node[cttl, anchor=west, align=left] at (0.1,-0.56) {cuPauliProp\\Distributed};
  \foreach \k in {0,...,35}
  {
    \node[draw=black!60, line width=0.3pt, fill=black!22, anchor=west, minimum width=0.2cm, minimum height=0.4cm, inner sep=0pt] at ({2.3+0.314*\k},-0.4) {};
    \fill[dred] ({2.3+0.314*\k+0.218},-0.12) rectangle ++(0.065,-0.56);
  }
  \node[inote, anchor=west] at (2.25,-0.9) {One Gate};
  \fill[dred] (3.75,-0.8) rectangle ++(0.065,-0.2);
  \node[inote, anchor=west] at (3.87,-0.9) {Global Merge on All GPUs};
  \node[key, anchor=east] at (13.8,-0.9) {22 Collective Calls per Gate};
  % ================= SectorQ: four nodes, three windows =================
  \node[frame, anchor=north west, minimum width=13.9cm, minimum height=3.3cm] at (0,-1.3) {};
  \node[lttl, anchor=north west] at (0.1,-1.33) {SectorQ};
  \foreach \x/\t in {4.05/Window 8, 7.95/Window 9, 11.85/Window 10}
    \node[cttl] at (\x,-1.92) {\t: 127 Gates};
  \foreach \n/\f/\a/\b/\c/\d/\e/\g in {0/sa/0/4/8/12/16/20, 1/sb/1/5/9/13/17/21, 2/sc/2/6/10/14/18/22, 3/sd/3/7/11/15/19/23}
  {
    \node[cttl, anchor=west] at (0.1,{-2.45-0.55*\n}) {Node \n};
    \foreach \x in {2.3, 6.2, 10.1}
    {
      \node[gpu, anchor=west, minimum width=3.42cm, minimum height=0.44cm] at (\x,{-2.45-0.55*\n}) {};
      \foreach \i/\t in {0/\a, 1/\b, 2/\c, 3/\d, 4/\e, 5/\g}
        \node[sec, fill=\f] at ({\x+0.08+0.46*\i},{-2.45-0.55*\n}) {\t};
      \foreach \i/\ff in {0/sa, 1/sb, 2/sc, 3/sd}
        \node[pc, fill=\ff] at ({\x+2.92+0.11*\i},{-2.45-0.55*\n}) {};
    }
  }
  % the exchange at the two boundaries
  \foreach \x in {5.8, 9.7}
  {
    \fill[dblue!25] (\x,-2.2) rectangle ++(0.4,-2.15);
    \draw[dfat, dblue] ({\x+0.2},-2.27) -- ({\x+0.2},-4.28);
  }
  \node[inote, text=dblue] at (6.0,-1.92) {MPI};
  \node[inote, text=dblue] at (9.9,-1.92) {MPI};
  \node[key, anchor=north east] at (13.8,-1.36) {No MPI Inside a Window};
  \node[pc, fill=sa] at (2.3,-1.52) {}; \node[pc, fill=sb] at (2.41,-1.52) {}; \node[pc, fill=sc] at (2.52,-1.52) {}; \node[pc, fill=sd] at (2.63,-1.52) {};
  \node[inote, anchor=west] at (2.78,-1.52) {Output, Divided for the Next Window};
  % ================= measured number of terms =================
  \node[layer, anchor=north west, minimum width=13.9cm, minimum height=1.82cm] at (0,-4.78) {};
  \node[cttl, anchor=west, align=left] at (0.1,-5.5) {Terms\\(Measured)};
  \foreach \x in {5.8, 9.7}
    \fill[dblue!25] (\x,-4.9) rectangle ++(0.4,-1.38);
  \fill[black!28] \curvepts -- (13.6,-6.22) -- (2.3,-6.22) -- cycle;
  \draw[black!85, line width=0.7pt] plot coordinates {\curvepts};
  \draw[black!70, line width=0.5pt] (2.3,-6.22) -- (13.6,-6.22);
  \foreach \x in \truncticks
    \draw[black!70, line width=0.5pt] (\x,-6.22) -- (\x,-6.31);
  \node[key, anchor=west, inner sep=1pt] at (4.3,-5.08) {1.36 Billion at the Peak};
  \node[key, anchor=south, inner sep=1pt] at (6.0,-5.78) {159M};
  \node[inote, anchor=south, inner sep=1pt] at (9.9,-5.82) {147M};
  \node[inote, anchor=west] at (10.5,-5.08) {Ticks: Truncations};
\end{tikzpicture}
